\documentclass[11pt,a4paper]{article}

\usepackage[utf8]{inputenc}
\usepackage[T1]{fontenc}
\usepackage[spanish,es-noshorthands,es-nodecimaldot]{babel}
\usepackage[margin=2.5cm]{geometry}
\usepackage{graphicx}
\usepackage{float}
\usepackage{booktabs}
\usepackage{amsmath}
\usepackage{enumitem}
\usepackage{caption}
\usepackage[hidelinks]{hyperref}

% Las figuras y tablas se generan en ../output/
\graphicspath{{../output/}}

\setlist[itemize]{itemsep=2pt,topsep=4pt}
\captionsetup{font=small,labelfont=bf}

\begin{document}

% ---------------------------------------------------------------
% Portada
% ---------------------------------------------------------------
\begin{titlepage}
\centering
\vspace*{3cm}

{\LARGE\bfseries Trabajo Práctico N$^{\circ}$1\par}
\vspace{0.4cm}
{\large (Semestre de Otoño 2025)\par}
\vspace{0.3cm}
{\large\bfseries Big Data\par}

\vfill

{\large Profesora: María Noelia Romero\par}
\vspace{0.2cm}
{\large Tutor: Tomás Pacheco\par}
\vspace{0.2cm}
{\large Alumnos: Catalina Gallardo, Manuel Elorga\par}

\vfill
\end{titlepage}

% ===============================================================
\section*{Parte I: Familiarizándonos con la base EPH}
% ===============================================================

\subsection*{Ejercicio 1}

Para identificar si una persona está desocupada, el INDEC aplica tres criterios
fundamentales. El primero es no haber desarrollado ningún tipo de actividad económica
remunerada. El segundo es haber realizado acciones concretas para obtener trabajo en los
últimos 30 días, como enviar currículums o registrarse en agencias de empleo. El tercero es
estar disponible para incorporarse a un empleo en caso de recibir una oferta.

El INDEC diferencia tres grupos dentro del mercado de trabajo: la población ocupada, la
desocupada y la inactiva. La población ocupada está compuesta por aquellas personas que
han trabajado al menos una hora en la semana de referencia, ya sea de manera remunerada o
sin remuneración. La población inactiva está conformada por aquellas personas que no tienen
empleo y que no lo buscan de manera activa, como estudiantes, jubilados y personas dedicadas
a tareas del hogar. La población económicamente activa (PEA) incluye tanto a las personas
ocupadas como a las desocupadas.

La tasa de desocupación se calcula como el porcentaje de la población desocupada sobre la
población económicamente activa. Este indicador permite evaluar la evolución del mercado
laboral, ya que da a conocer la cantidad de personas que desean trabajar pero no encuentran
empleo.

El análisis de la desocupación es fundamental para comprender la dinámica del mercado de
trabajo y evaluar el impacto de las políticas económicas y sociales. Una tasa de desocupación
elevada puede reflejar problemas como la falta de oportunidades laborales en diferentes
sectores o crisis económicas que afectan la demanda de empleo.

\subsection*{Ejercicio 2}

\subsubsection*{Inciso b}

Para el heatmap de valores faltantes elegimos las siguientes diez variables:

\begin{itemize}
  \item \textbf{CH14}: último año de educación aprobado.
  \item \textbf{CH15\_COD}: código correspondiente al lugar de nacimiento (provincia o país extranjero).
  \item \textbf{CH16\_COD}: código correspondiente al lugar de residencia hace cinco años (provincia o país extranjero).
  \item \textbf{PP03C}: número de empleos o actividades laborales desempeñadas en la última semana.
  \item \textbf{PP03I}: búsqueda de trabajo adicional en los últimos 30 días.
  \item \textbf{PP04B\_COD}: código de la actividad económica de la empresa o institución en la que trabaja.
  \item \textbf{PP06H}: indica si la actividad es un negocio familiar.
  \item \textbf{PP07A}: tiempo continuo trabajado en el empleo actual.
  \item \textbf{PP10A}: tiempo de búsqueda de empleo en personas desocupadas.
  \item \textbf{PP11A}: sector (estatal, privado u otro) del último empleo de los desocupados con trabajo anterior.
\end{itemize}

\begin{figure}[H]
  \centering
  \includegraphics[width=0.85\textwidth]{{Heatmap 1}.png}
  \caption{Heatmap de valores faltantes, antes de la limpieza.}
  \label{fig:heatmap-pre}
\end{figure}

La Figura~\ref{fig:heatmap-pre} muestra que faltan bastantes datos sobre el último año
aprobado, tanto para 2004 como para 2024. Además, se observa una gran cantidad de valores
faltantes en las variables de código de lugar de nacimiento, de residencia y de actividad
económica. Por último, se reportan muchos valores faltantes en preguntas realizadas
únicamente a personas ocupadas o desocupadas en 2024. Esto no era un problema en 2004, lo
que lleva a pensar que puede haber diferencias en cómo se cargan los datos entre ambos
relevamientos.

\subsubsection*{Inciso c}

Eliminamos las siguientes observaciones sin sentido:

\begin{itemize}
  \item En la columna CH06 (edad) existían observaciones con el valor $-1$, que deberían
        representar ``menor de 1''. Se convirtieron a 0.
  \item En la columna CH14 (último año aprobado), en la base de 2004 aparecían valores
        incorrectos (98), los cuales fueron reemplazados por valores nulos. Para 2024, el
        valor 98 significa ``educación especial''.
  \item En la columna PP06C (ingresos) se encontraron valores negativos distintos de $-9$
        (que indica ``no sabe / no contesta''). Los valores $-7$ y $-8$ fueron reemplazados
        por valores nulos.
  \item En la columna PP06D (ingresos), al igual que en PP06C, se reemplazaron los valores
        negativos distintos de $-9$, convirtiendo $-7$ y $-8$ en valores nulos.
\end{itemize}

Fuera de esos reemplazos, se verificó que todas las columnas tuvieran sus valores dentro del
rango de respuestas válidas. Entre otros controles, se verificó que la cantidad de horas
trabajadas en la última semana no excediera las 168 horas que tiene una semana.

\subsubsection*{Inciso d}

Para limpiar la base de datos se realizaron los siguientes cambios:

\begin{itemize}
  \item Se eliminaron las columnas que existían solo en una de las bases fusionadas, para
        garantizar uniformidad.
  \item Se transformó la variable MAS\_500 en binaria, asignando 1 si el aglomerado tiene
        más de 500.000 habitantes y 0 en caso contrario.
  \item Se estableció que si la persona es menor de dos años (CH09 = 3) o nunca asistió a un
        establecimiento educativo (CH10 = 3), entonces no aprobó ningún nivel educativo
        (CH14 = 0), reemplazando así los valores nulos.
  \item Se ajustaron las variables CH15\_COD y CH16\_COD, que solo deben tomar valores si
        CH15 o CH16 son 3, 4 o 5, asignándoles 0 en los demás casos (que interpretamos como
        ``pregunta no aplica'').
  \item Se ajustó PP04B\_COD para que solo tome valores si PP04A es 3, reemplazando los
        valores nulos por 0.
  \item Se ajustó PP06H para que solo tome valores si PP06E es 2 o 3, reemplazando los
        valores nulos por 0.
  \item Se convirtieron en variables binarias PP05F, asignando 1 si el negocio trabaja con
        un solo cliente y 2 si trabaja con múltiples clientes, y PP07F5 y PP07G\_59,
        asignando 1 si la respuesta es afirmativa.
  \item Se reemplazaron por 0 los valores faltantes en variables que solo aplican a personas
        ocupadas o desocupadas, cuando el individuo no pertenece a esa categoría, para
        disminuir la cantidad de valores nulos.
  \item Se eliminaron variables no necesarias: IMPUTA (sin observaciones), PP09A,
        PP09A\_ESP, PP09B, PP09C y PP09C\_ESP (dependientes de otras variables o aplicables
        solo a ciertas regiones).
  \item Se convirtieron a enteros las variables que deberían tener valores discretos,
        excluyendo los montos de ingreso y la variable CODUSU (el código que distingue
        viviendas).
  \item Se garantizó que las variables de ingreso fueran de tipo \textit{float}, para
        preservar precisión en los cálculos.
\end{itemize}

\subsubsection*{Inciso e}

Luego de la limpieza, el heatmap de valores faltantes se ve como muestra la
Figura~\ref{fig:heatmap-post}.

\begin{figure}[H]
  \centering
  \includegraphics[width=0.85\textwidth]{{Heatmap 2}.png}
  \caption{Heatmap de valores faltantes, después de la limpieza.}
  \label{fig:heatmap-post}
\end{figure}

En la Figura~\ref{fig:heatmap-post} vemos que con la limpieza se logró reemplazar casi todos
los valores faltantes, con la excepción de aquellos de la columna CH14 (último año escolar
aprobado). Si bien se hizo inferencia respecto de algunas personas que no habían aprobado
ningún año ---como los menores de dos años, que aún no están en edad escolar, y aquellos que
reportaron nunca haber asistido a una institución educativa---, el resto de los casos no
admitía una imputación razonable y se conservaron como faltantes.

% ===============================================================
\section*{Parte II: Primer análisis exploratorio}
% ===============================================================

\subsection*{Ejercicio 3}

La Figura~\ref{fig:sexo} ilustra la comparación de la cantidad de hombres y mujeres en 2004 y
2024, lo que permite observar cómo evolucionó la composición por sexo a lo largo del tiempo.

En 2004 se registraron aproximadamente 2291 hombres y 2530 mujeres, es decir, una diferencia
de 239 personas a favor de las mujeres. En 2024 se observan alrededor de 2143 hombres y 2370
mujeres, lo que da nuevamente una brecha a favor de las mujeres, esta vez de 227 personas.
Cabe destacar que se presenta una leve disminución respecto de la diferencia previa. Además,
tanto en hombres como en mujeres hay una reducción de la cantidad de observaciones entre 2004
y 2024. En definitiva, la composición por sexo experimenta variaciones relativamente
pequeñas en el período.

\begin{figure}[H]
  \centering
  \includegraphics[width=0.75\textwidth]{{Grafico barras 1}.png}
  \caption{Composición por sexo en 2004 y 2024.}
  \label{fig:sexo}
\end{figure}

\subsection*{Ejercicio 4}

Las Figuras~\ref{fig:matriz2004} y~\ref{fig:matriz2024} presentan las matrices de correlación
para 2004 y 2024 entre las variables CH04 (sexo), CH06 (edad), CH07 (estado conyugal), CH08
(cobertura médica), NIVEL\_ED (nivel educativo), ESTADO (ocupado, desocupado, inactivo),
CAT\_INAC (categoría de inactividad) e IPCF (ingreso per cápita familiar).

En ambos años, CH04 (sexo) muestra correlaciones bajas con el resto de las variables, lo que
indica que ser hombre o mujer no tiene una asociación significativa con factores como la
edad, el nivel educativo o el ingreso familiar. Asimismo, CH07 (estado conyugal) mantiene
correlaciones débiles con las otras variables, lo que sugiere que el estado civil no está
fuertemente vinculado con la educación o el empleo. Por su parte, la relación entre CH06
(edad) y CAT\_INAC sigue siendo positiva, lo que confirma que a medida que aumenta la edad
también lo hace la probabilidad de ser inactivo (jubilado, estudiante, etc.).

En relación con CH08 (cobertura médica), en ambas matrices se observa una correlación
positiva con NIVEL\_ED y con IPCF, lo que indica que a mayor nivel educativo y mayores
ingresos es más probable contar con cobertura médica. En 2024 esta relación se intensifica, lo
que sugiere que el acceso a la salud depende cada vez más de factores económicos y
educativos. Es importante señalar que la correlación negativa que se observa en la matriz no
implica una relación inversa, sino que depende de la codificación de la variable, donde
valores más bajos indican una mejor cobertura. Además, la relación entre NIVEL\_ED e IPCF,
aunque sigue siendo positiva, disminuyó en 2024, lo que refleja posibles cambios en la
estructura del mercado laboral o en la accesibilidad de los ingresos para los niveles
educativos más altos.

Por otro lado, ESTADO (condición de actividad) mantiene una correlación negativa con
CAT\_INAC, ya que una persona no puede ser activa e inactiva al mismo tiempo, y esta relación
se mantiene estable en 2024. La correlación entre ESTADO e IPCF sigue siendo baja en ambos
años, lo que puede explicarse porque los hogares tienen otras fuentes de ingresos además del
empleo individual. Finalmente, CAT\_INAC se correlaciona negativamente tanto con ESTADO como
con CH06 (edad), lo que reafirma que la inactividad está estrechamente relacionada con la
edad y la condición laboral. En general, las matrices reflejan cambios en la dinámica laboral
y social en estas dos décadas, con una mayor dependencia del nivel educativo y de los
ingresos para el acceso a bienes y servicios como la cobertura médica.

\begin{figure}[H]
  \centering
  \includegraphics[width=0.72\textwidth]{{Matriz 2004}.png}
  \caption{Matriz de correlación, 2004.}
  \label{fig:matriz2004}
\end{figure}

\begin{figure}[H]
  \centering
  \includegraphics[width=0.72\textwidth]{{Matriz 2024}.png}
  \caption{Matriz de correlación, 2024.}
  \label{fig:matriz2024}
\end{figure}

% ===============================================================
\section*{Parte III: Conociendo a los ocupados y desocupados}
% ===============================================================

\subsection*{Ejercicio 5}

Hay 323 desocupados y 5090 inactivos en la muestra. La media del ingreso per cápita familiar
es de 96894.0 para los ocupados, 33435.31 para los desocupados y 66566.35 para los inactivos.

\subsection*{Ejercicio 6}

Hay 6 personas en la muestra que no respondieron sobre su condición de actividad, todas
en 2024.

\subsection*{Ejercicio 7}

La Figura~\ref{fig:pea} muestra la distribución entre población económicamente activa y no
activa. Podemos ver que hay más personas no activas que activas, lo que puede deberse a la
gran cantidad de niños presentes en la muestra. Se aprecia también una disminución de la
cantidad de observaciones entre 2004 y 2024.

\begin{figure}[H]
  \centering
  \includegraphics[width=0.75\textwidth]{{Grafico barras 2}.png}
  \caption{Composición por PEA, 2004 y 2024.}
  \label{fig:pea}
\end{figure}

\subsection*{Ejercicio 8}

La Figura~\ref{fig:pet} muestra la distribución entre personas en edad laboral y fuera de la
edad laboral. Vemos que la mayor parte de la muestra, en ambos años, está en edad laboral.
Sin embargo, esto no significa que efectivamente estén ocupados ni que estén buscando
trabajo. Comparando con la figura anterior, muchas personas en edad laboral no están buscando
trabajo, lo que puede deberse a que sean estudiantes o a que pertenezcan a otra categoría de
inactividad.

\begin{figure}[H]
  \centering
  \includegraphics[width=0.75\textwidth]{{Grafico barras 3}.png}
  \caption{Composición por PET, 2004 y 2024.}
  \label{fig:pet}
\end{figure}

\subsection*{Ejercicio 9}

En la muestra, la desocupación bajó en 2024 respecto de 2004: hubo 223 desocupados en 2004 y
100 en 2024.

\subsubsection*{Inciso a}

La Figura~\ref{fig:educacion} muestra las proporciones de desocupados por nivel educativo en
2004 y 2024, y permite observar cambios importantes en la relación entre el nivel educativo y
la desocupación a lo largo del tiempo. Para calcular estas proporciones se contó a los
desocupados de cada grupo y se dividió por la población económicamente activa del mismo
grupo, de modo de dejar de lado a las personas inactivas. Un ejemplo de esto serían los niños
que no terminaron la primaria porque están aún en edad escolar.

Los principales resultados son los siguientes:

\begin{itemize}
  \item Entre las personas económicamente activas con nivel educativo \textbf{primaria
        incompleta}, el 15.25\% estaba desocupado en 2024, frente al 3.61\% en 2004.
  \item Con \textbf{primaria completa}, el 1.76\% en 2024, frente al 8.08\% en 2004.
  \item Con \textbf{secundaria incompleta}, el 6.26\% en 2024, frente al 13.79\% en 2004.
  \item Con \textbf{secundaria completa}, el 4.59\% en 2024, frente al 11.09\% en 2004.
  \item Con \textbf{universitaria incompleta}, el 5.76\% en 2024, frente al 15.86\% en 2004.
  \item Con \textbf{universitaria completa}, el 2.89\% en 2024, frente al 7.58\% en 2004.
  \item \textbf{Sin instrucción}, el 0.0\% en 2024, frente al 20.0\% en 2004. Esta gran
        reducción se explica principalmente por el hecho de que en 2024 solo dos personas de
        la muestra reportaron nivel educativo ``sin instrucción''.
\end{itemize}

En general, la desocupación disminuyó entre 2004 y 2024 en casi todos los niveles educativos,
lo que indica un mejor panorama laboral. La única excepción es primaria incompleta, donde la
proporción de desocupados aumentó, posiblemente por mayores exigencias del mercado laboral
para los niveles de calificación más bajos. La mejora fue más notable en los niveles más
altos de educación (universitaria completa e incompleta). Este análisis destaca la importancia
de la educación en la inserción laboral, con una clara ventaja para quienes poseen niveles
educativos más altos. Sin embargo, es importante considerar el reducido tamaño de la muestra
en cada nivel.

\begin{figure}[H]
  \centering
  \includegraphics[width=0.85\textwidth]{{Grafico barras 4}.png}
  \caption{Proporción de desocupados por nivel educativo, 2004 y 2024.}
  \label{fig:educacion}
\end{figure}

\subsubsection*{Inciso b}

La Figura~\ref{fig:edad} compara las tasas de desocupación por grupo etario en 2004 y 2024, y
permite observar cómo la desocupación afecta a distintas edades y cómo cambiaron estas
proporciones a lo largo del tiempo.

\begin{itemize}
  \item \textbf{Grupos sin participación económica}: no hay personas económicamente activas
        en los grupos de 0 a 9 años ni de 90 a 99 años en ninguno de los dos años
        analizados, lo cual es consistente con lo esperado.

  \item \textbf{Desocupación en los jóvenes}: en 2004, los jóvenes de 10 a 19 y de 20 a 29
        años presentan las tasas más altas de desocupación, con valores de 19.05\% y
        19.19\% respectivamente. Esto podría deberse a la falta de experiencia laboral y a la
        transición de los estudios al empleo. En 2024 las tasas de estos grupos disminuyen
        considerablemente, a 5.88\% y 9.75\%. Aunque siguen siendo relativamente altas en
        comparación con otros grupos, la reducción es notable.

  \item \textbf{Desocupación en los adultos}: en 2004, los adultos de 30 a 49 años presentan
        tasas moderadas (7.23\% para 30 a 39 años y 7.02\% para 40 a 49 años), lo que refleja
        una etapa de mayor estabilidad laboral que la de los grupos más jóvenes. En 2024
        estas tasas se reducen a 3.79\% y 3.68\% respectivamente, lo que indica un panorama
        más favorable para los adultos en edad productiva.

  \item \textbf{Desocupación en los adultos mayores}: en 2004, los grupos de 50 a 59, 60 a 69
        y 70 a 79 años tienen tasas relativamente bajas: 5.52\%, 5.15\% y 4.76\%
        respectivamente. En 2024 estas tasas disminuyen a 2.05\% y 2.25\% en los dos primeros
        grupos, mientras que en el de 70 a 79 años sube levemente a 4.35\%. Esto podría
        deberse a una menor participación activa de estos grupos en el mercado laboral.
        Para ambos años no hay desocupados en el grupo de 80 a 89, aunque sí hay personas
        económicamente activas, probablemente porque a esa edad quienes siguen en el mercado
        laboral lo hacen por decisión propia y podrían pasar a ser inactivos si lo desearan.
\end{itemize}

\begin{figure}[H]
  \centering
  \includegraphics[width=0.85\textwidth]{{Grafico barras 5}.png}
  \caption{Proporción de desocupados por edad, 2004 y 2024.}
  \label{fig:edad}
\end{figure}

\subsubsection*{Inciso c}

La Figura~\ref{fig:conyugal} presenta las tasas de desocupación por estado civil entre 2004 y
2024, que revelan diferencias significativas según la categoría y un descenso generalizado de
la desocupación en el período.

\begin{itemize}
  \item \textbf{Unido/a}: 3.77\% de desocupados entre los económicamente activos en 2024,
        frente al 7.89\% en 2004.
  \item \textbf{Casado/a}: 2.31\% en 2024, frente al 5.59\% en 2004.
  \item \textbf{Separado/a o divorciado/a}: 3.49\% en 2024, frente al 10.06\% en 2004.
  \item \textbf{Viudo/a}: 6.82\% en 2024, frente al 5.66\% en 2004.
  \item \textbf{Soltero/a}: 7.46\% en 2024, frente al 19.08\% en 2004.
\end{itemize}

Se puede apreciar que la disminución de la desocupación es notable en todas las categorías
entre 2004 y 2024, salvo un leve aumento en la categoría viudo/a. Además, las personas
solteras y separadas o divorciadas son más vulnerables a la desocupación en ambos años,
posiblemente debido a factores como una menor estabilidad económica o social. Por último, las
personas casadas presentan las tasas más bajas de desocupación en ambos períodos, lo que
puede estar asociado a mayores redes de apoyo o a una mayor estabilidad laboral.

\begin{figure}[H]
  \centering
  \includegraphics[width=0.85\textwidth]{{Grafico barras 6}.png}
  \caption{Proporción de desocupados por estado conyugal, 2004 y 2024.}
  \label{fig:conyugal}
\end{figure}

\subsection*{Ejercicio 10}

La Figura~\ref{fig:tasas} muestra la comparación entre dos tipos de tasas de desocupación
para 2004 y 2024: la calculada según la definición del INDEC y una tasa alternativa basada en
la economía laboral.

\paragraph{Tasa de desocupación del INDEC.}
En 2004 fue del 5.26\%, mientras que en 2024 disminuyó al 2.36\%. Esto refleja una reducción
significativa en el porcentaje de desocupados en relación con la población económicamente
activa (PEA), lo que podría estar vinculado a una mejora en las condiciones laborales o a
cambios en la estructura económica.

\paragraph{Tasa de desocupación alternativa.}
En 2004 fue del 3.65\%, y disminuyó al 1.64\% en 2024. Este cálculo utiliza como base la
población en edad de trabajar (PET), lo que implica un indicador más amplio, al no depender
exclusivamente de la decisión de participar en el mercado laboral.

La tasa del INDEC tiene varias ventajas. En primer lugar, se centra exclusivamente en la
población activa, lo que permite un indicador ajustado a las personas que efectivamente están
buscando empleo. Además, al ser una medida estándar utilizada internacionalmente, facilita
comparaciones entre distintas regiones o países, lo que la convierte en una herramienta
práctica para el análisis global del mercado laboral. Sin embargo, también presenta
desventajas: puede estar sesgada por la decisión endógena de las personas de participar en el
mercado laboral. Esto significa que, si la participación en la PEA disminuye, la tasa podría
aparentar una mejora aunque las condiciones laborales no necesariamente hayan mejorado.

Por otro lado, la tasa alternativa considera a toda la población en edad de trabajar,
eliminando el sesgo asociado a las decisiones individuales de participar o no en el mercado
laboral. Esto permite una visión más inclusiva y amplia de las condiciones laborales de la
economía. Su desventaja es que incluye a individuos que no buscan empleo por elección, lo que
podría dar una percepción inflada de la gravedad de la desocupación en ciertos contextos.

Ambas tasas tienen utilidad según el enfoque que se desee analizar, y combinarlas puede
proporcionar una evaluación más integral del mercado laboral. La diferencia entre ambas
sugiere que la medición alternativa proporciona una visión más conservadora del problema, ya
que reparte los desocupados sobre una base poblacional mayor y, por lo tanto, arroja niveles
sistemáticamente más bajos que la tasa oficial.

\begin{figure}[H]
  \centering
  \includegraphics[width=0.8\textwidth]{{Grafico barras 7}.png}
  \caption{Tasas de desocupación, 2004 y 2024.}
  \label{fig:tasas}
\end{figure}

\end{document}
